\chapter{Conclusões e Trabalhos Futuros} 
\label{cap:conclusao}

A automação da Alta Disponibilidade (HA) em bancos de dados relacionais representa um desafio crítico para a confiabilidade de sistemas de supervisão SCADA. Esta monografia propôs o desenvolvimento de uma arquitetura resiliente capaz de mitigar os impactos operacionais decorrentes de falhas de \textit{hardware} e isolamentos de rede, eliminando a dependência de intervenções manuais.

A implementação do agente orquestrador \textit{raft-ha-manager}, construído de forma nativa na linguagem Java e fundamentado no algoritmo de consenso Raft \cite{raft2014}, cumpriu o objetivo central da pesquisa. Os resultados demonstraram que o sistema reduz drasticamente o Tempo Médio de Recuperação (MTTR), que passa de horas, patamar comum em processos manuais de diagnóstico e clonagem, para a escala de segundos, garantindo o redirecionamento ininterrupto das operações do SCADA em cenários de queda do serviço.

Destaca-se como diferencial acadêmico a decisão arquitetural de não adotar soluções comerciais fechadas ou bibliotecas de consenso genéricas. O desenvolvimento do protocolo Raft a partir de seus fundamentos teóricos proporcionou um acoplamento fino ao ciclo de vida do PostgreSQL. Esse controle de baixo nível possibilitou a execução ativa de rotinas de \textit{Fencing} e \textit{Self-Healing}, além de viabilizar a renúncia voluntária da liderança do \textit{cluster} perante a degradação da saúde do banco local. Tais mecanismos asseguraram a consistência dos dados industriais e mitigaram o estado anômalo de \textit{Split-Brain}, comprovando o êxito quantitativo e estrutural do projeto.

\section{Trabalhos Futuros}

Embora a solução desenvolvida atenda aos requisitos de confiabilidade delineados no escopo desta monografia, as limitações teóricas e práticas observadas durante as baterias de testes abrem caminho para evoluções arquiteturais. Sugere-se o aprofundamento das seguintes temáticas como trabalhos futuros:

\begin{itemize}
    \item \textbf{Implementação da Fase de Pré-Votação (\textit{Pre-Vote}):} As anomalias de perturbação registradas no cenário de particionamento de rede comprovaram a vulnerabilidade do algoritmo clássico perante a elevação no escuro do mandato (\textit{term}) de nós isolados. A integração da extensão de \textit{Pre-Vote}, na qual a disponibilidade da maioria do \textit{cluster} deve ser atestada antes de qualquer incremento de mandato \cite{raft2014}, é a evolução imediata sugerida para impedir deposições indevidas de líderes sadios na reconexão de nós desatualizados.
    
    \item \textbf{Tolerância a Falhas Bizantinas (\verb|BFT| - \textit{Byzantine Fault Tolerant})\nomenclature{BFT}{Byzantine Fault Tolerant}:} O ecossistema industrial de redes SCADA é um alvo recorrente de ciberataques. O projeto atual baseia-se no modelo CFT (\textit{Crash Fault Tolerant}), que é incapaz de operar sob a presença de nós maliciosos. A exploração de protocolos baseados em Tolerância a Falhas Bizantinas, conforme estabelecido por \citeonline{Castro1999}, ou a incorporação de assinaturas criptográficas no tráfego gRPC, faz-se necessária para salvaguardar o consenso distribuído contra adulterações intencionais de votos ou mensagens.
    
    \item \textbf{Durabilidade Estrita com Invocação de \textit{fsync}:} Para transpor os limites da máquina virtual Java e garantir integridade absoluta contra quedas drásticas de energia, recomenda-se a substituição do método padrão de escoamento em memória por chamadas nativas de \textit{fsync} diretamente ao sistema operacional. Forçar a gravação física imediata dos bytes de log no dispositivo de armazenamento contorna o risco de perda de dados retidos no cache de página, solidificando a camada de durabilidade da máquina de estados do Raft.
    
    \item \textbf{Replicação Integrada do \textit{Write-Ahead Log} (WAL):} Como desafio avançado de engenharia de \textit{software}, sugere-se a centralização total do tráfego de dados. Ao invés de delegar a sincronização dos bancos de dados ao serviço de replicação contínua do PostgreSQL de forma paralela, propõe-se que versões futuras do agente interceptem e transportem o próprio arquivo de logs WAL do banco de dados através dos canais gRPC da máquina de estados do Raft. Tal arquitetura tornaria o protocolo responsável não apenas pela eleição e pelo tratamento de falhas, mas pela integridade de cada transação inserida no SCADA.
\end{itemize}
